% GENERATED FILE. Edit canonical lessons, then run npm run generate:books.
% canonical-source-hash: fnv1a64:c85ea66443f28d7a
% canonical-lessons: JA-C381-funatabi, JA-C381-kouro, JA-C381-kishi, JA-C381-kogu, JA-C381-tadayou

\chapter{The Ferry Crossing}
\label{ch:ja-the-ferry-crossing}
\hlchaptermodality{381}

\begin{chapteropening}
\textbf{By the end of this chapter:} \emph{I can talk about crossing the water by boat.}

Complete the last lesson of chapter 381: I can talk about crossing the water by boat.
\end{chapteropening}

\section[\texorpdfstring{\ja{ふなたび} (funatabi)}{ふなたび (funatabi)}]{\ja{ふなたび} (funatabi) — a journey by boat}

\label{lesson:JA-C381-funatabi}



\begin{quote}
Before the new one: say the Japanese for dozing off, then the Japanese for a mouthful.
\end{quote}

\subsection*{You'll want to know: \ja{ふなたび}}
\textbf{\ja{ふなたび}} — \emph{funatabi} — \textquotedblleft{}a journey by boat\textquotedblright{}.

\begin{quote}
\textbf{\ja{しまへの} \ja{ふなたびは} \ja{たのしかったです。}} — \emph{shima e no funatabi wa tanoshikatta desu} — \textquotedblleft{}The boat trip to the island was fun.\textquotedblright{}
\end{quote}

\begin{cousinweb}
\textbf{\ja{ふなたび}} (\emph{funatabi}) is a native Japanese word (wago), inherited rather than borrowed from Chinese, built from \emph{funa}, \textquotedblleft{}boat\textquotedblright{} (from \emph{fune}) and \emph{tabi}, \textquotedblleft{}journey\textquotedblright{}; its kanji were attached later as kun'yomi (JMdict).
\end{cousinweb}

\begin{grammarlens}[title={What you've built}]
One more word for getting around.
\end{grammarlens}

\subsection*{Your turn}
\begin{itemize}
\raggedright
  \item \emph{Say it:} \emph{funatabi}
  \item \emph{Say it:} \emph{funatabi}, once more
  \item \emph{Say it:} \emph{shima e no funatabi wa tanoshikatta desu} — The boat trip to the island was fun.
  \item \emph{Recall:} say the Japanese for dozing off, then the Japanese for a mouthful, then say \emph{funatabi} again
\end{itemize}

\begin{tcolorbox}[breakable,colback=teal!4,colframe=teal!35!black,title={Before you move on}]
What does \ja{ふなたび} mean? (\textquotedblleft{}A journey by boat\textquotedblright{}.) Say it once more.
\end{tcolorbox}
\section[\texorpdfstring{\ja{こうろ} (kōro)}{こうろ (kōro)}]{\ja{こうろ} (kōro) — a sea route; a course}

\label{lesson:JA-C381-kouro}



\begin{quote}
Before the new one: say the Japanese for a mouthful, then the Japanese for a journey by boat.
\end{quote}

\subsection*{You'll want to know: \ja{こうろ}}
\textbf{\ja{こうろ}} — \emph{kōro} — \textquotedblleft{}a sea route; a course\textquotedblright{}.

\begin{quote}
\textbf{\ja{この} \ja{こうろは} \ja{なみが} \ja{たかいです。}} — \emph{kono kōro wa nami ga takai desu} — \textquotedblleft{}The waves are high on this sea route.\textquotedblright{}
\end{quote}

\begin{cousinweb}
\textbf{\ja{こうろ}} (\emph{kōro}) is Sino-Japanese (kango): two Chinese characters read with their on'yomi, readings borrowed from Chinese. They are \emph{kō}, \textquotedblleft{}navigate\textquotedblright{} (Middle Chinese \emph{hang}) and \emph{ro}, \textquotedblleft{}route\textquotedblright{} (\emph{luH}) (KANJIDIC2; Middle Chinese in Baxter's transcription).
\end{cousinweb}

\begin{grammarlens}[title={What you've built}]
One more word for getting around.
\end{grammarlens}

\subsection*{Your turn}
\begin{itemize}
\raggedright
  \item \emph{Say it:} \emph{kōro}
  \item \emph{Say it:} \emph{kōro}, once more
  \item \emph{Say it:} \emph{kono kōro wa nami ga takai desu} — The waves are high on this sea route.
  \item \emph{Recall:} say the Japanese for a mouthful, then the Japanese for a journey by boat, then say \emph{kōro} again
\end{itemize}

\begin{tcolorbox}[breakable,colback=teal!4,colframe=teal!35!black,title={Before you move on}]
What does \ja{こうろ} mean? (\textquotedblleft{}A sea route; a course\textquotedblright{}.) Say it once more.
\end{tcolorbox}
\section[\texorpdfstring{\ja{きし} (kishi)}{きし (kishi)}]{\ja{きし} (kishi) — a shore; a bank}

\label{lesson:JA-C381-kishi}



\begin{quote}
Before the new one: say the Japanese for a journey by boat, then the Japanese for a sea route.
\end{quote}

\subsection*{You'll want to know: \ja{きし}}
\textbf{\ja{きし}} — \emph{kishi} — \textquotedblleft{}a shore; a bank\textquotedblright{}.

\begin{quote}
\textbf{\ja{ふねが} \ja{きしに} \ja{つきました。}} — \emph{fune ga kishi ni tsukimashita} — \textquotedblleft{}The boat reached the shore.\textquotedblright{}
\end{quote}

\begin{cousinweb}
\textbf{\ja{きし}} (\emph{kishi}) is a native Japanese word (wago), inherited rather than borrowed from Chinese; its kanji was attached later as a kun'yomi (JMdict).
\end{cousinweb}

\begin{grammarlens}[title={What you've built}]
One more word for getting around.
\end{grammarlens}

\subsection*{Your turn}
\begin{itemize}
\raggedright
  \item \emph{Say it:} \emph{kishi}
  \item \emph{Say it:} \emph{kishi}, once more
  \item \emph{Say it:} \emph{fune ga kishi ni tsukimashita} — The boat reached the shore.
  \item \emph{Recall:} say the Japanese for a journey by boat, then the Japanese for a sea route, then say \emph{kishi} again
\end{itemize}

\begin{tcolorbox}[breakable,colback=teal!4,colframe=teal!35!black,title={Before you move on}]
What does \ja{きし} mean? (\textquotedblleft{}A shore; a bank\textquotedblright{}.) Say it once more.
\end{tcolorbox}
\section[\texorpdfstring{\ja{こぐ} (kogu)}{こぐ (kogu)}]{\ja{こぐ} (kogu) — to row; to pedal}

\label{lesson:JA-C381-kogu}



\begin{quote}
Before the new one: say the Japanese for a sea route, then the Japanese for a shore.
\end{quote}

\subsection*{You'll want to know: \ja{こぐ}}
\textbf{\ja{こぐ}} — \emph{kogu} — \textquotedblleft{}to row; to pedal\textquotedblright{}.

\begin{quote}
\textbf{\ja{ちちと} \ja{ふねを} \ja{こぎました。}} — \emph{chichi to fune o kogimashita} — \textquotedblleft{}I rowed a boat with my father.\textquotedblright{}
\end{quote}

Polite \textbf{\ja{こぎます}}, \emph{kogimasu}; a godan verb. Past (polite): \textbf{\ja{ちちが} \ja{ふねを} \ja{こぎました。}}, \emph{chichi ga fune o kogimashita}.

\begin{cousinweb}
\textbf{\ja{こぐ}} (\emph{kogu}) is a native Japanese word (wago), inherited rather than borrowed from Chinese; its kanji was attached later as a kun'yomi (JMdict).
\end{cousinweb}

\begin{grammarlens}[title={What you've built}]
One more word for getting around.
\end{grammarlens}

\subsection*{Your turn}
\begin{itemize}
\raggedright
  \item \emph{Say it:} \emph{kogu}
  \item \emph{Say it:} \emph{kogu}, once more
  \item \emph{Say it:} \emph{chichi to fune o kogimashita} — I rowed a boat with my father.
  \item \emph{Recall:} say the Japanese for a sea route, then the Japanese for a shore, then say \emph{kogu} again
\end{itemize}

\begin{tcolorbox}[breakable,colback=teal!4,colframe=teal!35!black,title={Before you move on}]
What does \ja{こぐ} mean? (\textquotedblleft{}To row; to pedal\textquotedblright{}.) Say it once more.
\end{tcolorbox}
\section[\texorpdfstring{\ja{ただよう} (tadayou)}{ただよう (tadayou)}]{\ja{ただよう} (tadayou) — to drift; to float}

\label{lesson:JA-C381-tadayou}



\begin{quote}
Before the new one: say the Japanese for a shore, then the Japanese for to row.
\end{quote}

\subsection*{You'll want to know: \ja{ただよう}}
\textbf{\ja{ただよう}} — \emph{tadayou} — \textquotedblleft{}to drift; to float\textquotedblright{}.

\begin{quote}
\textbf{\ja{ちいさい} \ja{ふねが} \ja{うみを} \ja{ただよっています。}} — \emph{chiisai fune ga umi o tadayotte imasu} — \textquotedblleft{}A small boat is drifting on the sea.\textquotedblright{}
\end{quote}

Polite \textbf{\ja{ただよいます}}, \emph{tadayoimasu}; a godan verb. Past (polite): \textbf{\ja{ふねが} \ja{うみを} \ja{ただよいました。}}, \emph{fune ga umi o tadayoimashita}.

\begin{cousinweb}
\textbf{\ja{ただよう}} (\emph{tadayou}) is a native Japanese word (wago), inherited rather than borrowed from Chinese; its kanji was attached later as a kun'yomi (JMdict).
\end{cousinweb}

\begin{grammarlens}[title={What you've built}]
One more word for getting around.
\end{grammarlens}

\subsection*{Your turn}
\begin{itemize}
\raggedright
  \item \emph{Say it:} \emph{tadayou}
  \item \emph{Say it:} \emph{tadayou}, once more
  \item \emph{Say it:} \emph{chiisai fune ga umi o tadayotte imasu} — A small boat is drifting on the sea.
  \item \emph{Recall:} say the Japanese for a shore, then the Japanese for to row, then say \emph{tadayou} again
\end{itemize}

\begin{tcolorbox}[breakable,colback=teal!4,colframe=teal!35!black,title={Before you move on}]
What does \ja{ただよう} mean? (\textquotedblleft{}To drift; to float\textquotedblright{}.) Say it once more.
\end{tcolorbox}
